% Fragment -- inclus de main.tex prin \input
\clearpage
\setpartlabel{themeCss}{\faPaintBrush}{Faza 2: CSS}
\setcounter{section}{0}
\phantomsection
\addcontentsline{toc}{part}{\protect\textbf{Faza 2 \textendash{} CSS}}
\handout{Faza 2 \textendash{} CSS}{Stilizare. Layout. Responsivitate.}

CSS (\emph{Cascading Style Sheets}) controleaza aspectul vizual al documentelor HTML. Numele provine de la \emph{cascada}: cand mai multe reguli vizeaza acelasi element, browserul aplica un algoritm de prioritate (cascada) pentru a alege valoarea finala.

Capitolul prezinta selectorii, modelul de cutie, unitatile de masura, variabilele CSS, sistemele de layout flexbox si grid, abordarea \emph{mobile-first}, si suportul nativ pentru tema dark/light.

\bigskip

\begin{outcomes}
Dupa capitolul asta vei \textcommabelow{s}ti:
\begin{itemize}[itemsep=1pt]
  \item Folose\textcommabelow{s}te selectorii CSS \textcommabelow{s}i in\textcommabelow{t}elege specificitatea ca triplet \code{(id, clasa, tag)}.
  \item Aplica modelul de cutie cu \code{box-sizing: border-box} pentru calcule predictibile.
  \item Define\textcommabelow{s}te variabile CSS in \code{:root} si le refolose\textcommabelow{s}te prin \code{var()}.
  \item Construie\textcommabelow{s}te layout-uri cu flexbox (6 proprieta\textcommabelow{t}i cheie).
  \item Scrie reguli mobile-first cu \code{@media (min-width: ...)}.
  \item Pozitioneaza elemente cu \code{position} \textcommabelow{s}i creeaza header sticky.
  \item Aplica anima\textcommabelow{t}ii cu \code{transition} si \code{@keyframes}, respectand \code{prefers-reduced-motion}.
  \item Combina \code{transform}, \code{transition} si \code{cubic-bezier()} pentru micro-interactiuni profesionale (hover-lift, fade-in, spinner).
  \item \faRocket\ \emph{Bonus self-study:} transformari 3D (flip-card), \code{animation-timeline} pentru scroll-driven, \code{filter}/\code{backdrop-filter}, \code{clip-path}, container queries.
\end{itemize}
\end{outcomes}

\begin{whymatters}
CSS face diferen\textcommabelow{t}a intre pagina-document \textcommabelow{s}i pagina-produs. Aceea\textcommabelow{s}i structura HTML poate parea neprofesionala sau profesionala, in func\textcommabelow{t}ie de stilizare. Mai mult: CSS-ul prost arhitecturat (cu \code{!important}, specificitate ridicata, valori absolute) genereaza ore de depanare. Conventiile prezentate aici (variabile, mobile-first, flexbox) sunt cele aplicate in practica industriala.
\end{whymatters}

\begin{nota}[Pre-rechizite]
Capitolul presupune in\textcommabelow{t}elegerea tag-urilor HTML semantice \textcommabelow{s}i a structurii \code{<head>}/\code{<body>} (Faza 1).
\end{nota}

% =========================================================
\section{Conectarea fisierului CSS la HTML}

CSS poate fi inclus in trei moduri. Conventia recomandata este externalizarea intr-un fisier separat:

\begin{lstlisting}[style=html, caption={Stilurile externe (recomandat)}]
<head>
  <link rel="stylesheet" href="style.css">
</head>
\end{lstlisting}

\begin{dano}[Asa nu \textendash{} stiluri inline imprastiate]
\begin{lstlisting}[style=html, numbers=none]
<button style="background: blue; color: white; padding: 10px;">Click</button>
\end{lstlisting}
\textbf{Probleme.} Imposibil de reutilizat; stilul nu poate fi schimbat global; specificitatea ridicata creeaza conflicte; pagina HTML devine ilizibila.
\end{dano}

\begin{dado}[Asa da \textendash{} clasa cu stil extern]
\begin{lstlisting}[style=html, numbers=none]
<button class="cta">Click</button>
\end{lstlisting}
In \code{style.css}: \code{.cta \{ background: blue; color: white; padding: 10px; \}}
\end{dado}

% =========================================================
\section{Selectori}

Selectorii determina elementele asupra carora se aplica un set de reguli. Categoriile principale:

\begin{lstlisting}[style=css, caption={Selectori uzuali}]
/* Selector de tip: vizeaza toate elementele <h1> */
h1 { color: navy; }

/* Selector de clasa: reutilizabil, prefixat cu . */
.cta { background: orange; }

/* Selector de id: unic pe pagina, prefixat cu # */
#contact { padding: 2rem; }

/* Selector de atribut */
input[type="email"] { border-color: blue; }

/* Combinator descendent: <a> oriunde sub <nav> */
nav a { text-decoration: none; }

/* Combinator copil direct: <li> direct sub <ul> */
ul > li { margin-bottom: 0.5rem; }

/* Combinator frate adiacent: <p> imediat dupa <h2> */
h2 + p { margin-top: 0; }

/* Pseudo-clase: stari ale elementului */
.cta:hover { background: red; }
.cta:focus-visible { outline: 2px solid blue; }
.cta:disabled { opacity: 0.5; }

/* Pseudo-elemente: parti generate */
p::first-letter { font-size: 2em; }
\end{lstlisting}

\subsection{Specificitate \textendash{} cum se rezolva conflictele}

Cand mai multe reguli vizeaza acelasi element, browserul calculeaza o specificitate pentru fiecare si o aplica pe cea cu valoare maxima. Calculul foloseste un triplet \code{(a, b, c)}:

\begin{itemize}[itemsep=1pt]
  \item \code{a} \textendash{} numarul de id-uri.
  \item \code{b} \textendash{} numarul de clase, atribute, pseudo-clase.
  \item \code{c} \textendash{} numarul de tag-uri si pseudo-elemente.
\end{itemize}

Comparatia se face lexicografic: \code{(0,1,0)} pierde fata de \code{(1,0,0)} indiferent de \code{b} si \code{c}. Stilurile inline au specificitate \code{(1,0,0,0)}, peste orice selector. Cuvantul cheie \code{!important} suprascrie tot, dar este considerat un \emph{antipattern} \textendash{} indica faptul ca arhitectura CSS este compromisa.

\begin{dano}[Asa nu \textendash{} folosirea \texttt{!important} ca solutie]
\code{.btn \{ background: red !important; \}}\\
\textbf{Problema.} Cand alta regula are nevoie sa schimbe culoarea, va trebui si ea sa foloseasca \code{!important}, escaladand conflictul.
\end{dano}

\begin{dado}[Asa da \textendash{} specificitate proportionala]
\code{.btn \{ background: blue; \} .btn.btn-danger \{ background: red; \}}\\
\textbf{Solutie:} clase modificator (\code{.btn-danger}) ridica specificitatea natural cu \code{(0,2,0)} fata de \code{(0,1,0)}.
\end{dado}

% =========================================================
\section{Modelul de cutie}

Fiecare element redat de browser este o cutie compusa din patru straturi concentrice: \emph{content}, \emph{padding}, \emph{border}, \emph{margin}.

\boxmodelfigure

\begin{itemize}
  \item \textbf{content} \textendash{} zona ocupata de continutul efectiv (text, imagine).
  \item \textbf{padding} \textendash{} spatiul dintre continut si chenar.
  \item \textbf{border} \textendash{} chenarul.
  \item \textbf{margin} \textendash{} spatiul exterior, intre cutia curenta si elementele adiacente.
\end{itemize}

\subsection{Implicit vs \texttt{border-box}}

Comportamentul implicit (\code{box-sizing: content-box}) include in \code{width} doar continutul. Adaugarea de \code{padding} sau \code{border} mareste latimea totala. Acest comportament face calculele predispuse la erori.

\begin{dano}[Asa nu \textendash{} content-box, latimi imprevizibile]
\begin{lstlisting}[style=css, numbers=none]
.card {
  width: 300px;
  padding: 20px;
  border: 1px solid #ccc;
  /* Latime totala: 300 + 20*2 + 1*2 = 342px (oversize) */
}
\end{lstlisting}
\end{dano}

\begin{dado}[Asa da \textendash{} border-box global, latimi predictibile]
\begin{lstlisting}[style=css, numbers=none]
*, *::before, *::after { box-sizing: border-box; }

.card {
  width: 300px;
  padding: 20px;
  border: 1px solid #ccc;
  /* Latime totala: 300px (padding si border absorbite in width) */
}
\end{lstlisting}
\end{dado}

% =========================================================
\section{Unitati de masura}

CSS suporta zeci de unitati. In practica, doar 5\textendash{}6 acopera majoritatea cazurilor:

\begin{itemize}
  \item \code{px} \textendash{} pixel absolut. Pentru border-uri si dimensiuni de detaliu.
  \item \code{rem} \textendash{} relativ la font-size-ul \code{<html>} (implicit 16px). Pentru font-size si spatiere uniforma.
  \item \code{em} \textendash{} relativ la font-size-ul elementului parinte. Util pentru padding-uri scaland cu textul.
  \item \code{\%} \textendash{} relativ la dimensiunea parintelui pe acea axa.
  \item \code{vw} / \code{vh} \textendash{} 1\% din latimea / inaltimea viewport-ului.
  \item \code{ch} \textendash{} latimea unui caracter ,,0''. Util pentru \code{max-width} pe paragrafe.
\end{itemize}

\begin{ponturi}
Pentru font-size si spatiere generala, \code{rem} este alegerea recomandata. Cand utilizatorul mareste font-size-ul implicit din browser (Settings $\to$ Appearance), tot site-ul scaleaza proportional. \code{px} pe font-size ignora aceasta preferinta.
\end{ponturi}

\begin{lstlisting}[style=css, caption={Compunere uzuala}]
:root {
  font-size: 100%;          /* echivalent 16px in browserele standard */
}

body { font-size: 1rem; }   /* 16px */
h1   { font-size: 2.5rem; } /* 40px */
.container { max-width: 64rem; }    /* 1024px */

p { max-width: 65ch; }      /* lungime optima de citire */

.hero {
  min-height: 80vh;         /* 80% din inaltimea ferestrei */
  padding: 2rem;
}
\end{lstlisting}

% =========================================================
\section{Variabile CSS (custom properties)}

Variabilele CSS permit definirea valorilor o singura data si reutilizarea lor pe parcursul intregului document. Modificarea unei singure declaratii in \code{:root} se propaga automat in toate locurile in care variabila este folosita.

\begin{lstlisting}[style=css]
:root {
  /* Paleta */
  --primary: #0673BA;
  --accent:  #FAA51E;
  --bg:      #FFFCF6;
  --text:    #1A1A1A;

  /* Spatiere consistenta */
  --space-1: 0.5rem;
  --space-2: 1rem;
  --space-3: 2rem;

  /* Layout */
  --max-width: 64rem;
  --radius: 0.5rem;
}

body {
  background: var(--bg);
  color: var(--text);
}

h1 { color: var(--primary); }
.cta { background: var(--accent); }
\end{lstlisting}

Variabilele CSS sunt \emph{scoped}: redefinirea intr-un selector descendent suprascrie valoarea pentru subarbore. Acest comportament permite teme contextuale fara duplicarea regulilor:

\begin{lstlisting}[style=css, caption={Tema schimbata local}]
.card-dark {
  --bg:   #1F2937;
  --text: #F1F5F9;
  /* toate copiii care folosesc var(--bg) si var(--text) preiau noile valori */
}
\end{lstlisting}

% =========================================================
\section{Proprietatea \texttt{display} \textendash{} valori}

\code{display} controleaza modul fundamental in care elementul se aseaza in pagina. Cele 6 valori pe care le intalnesti in 99\% din cazuri:

\begin{tcolorbox}[
  enhanced, breakable, sharp corners, arc=0pt,
  colback=accentSoft, colframe=partcolor, boxrule=0pt, leftrule=3pt,
  left=12pt, right=12pt, top=8pt, bottom=8pt, fontupper=\small,
]
\renewcommand{\arraystretch}{1.4}
\begin{tabular}{@{}p{0.22\linewidth}@{\hspace{8pt}}p{0.74\linewidth}@{}}
{\sffamily\bfseries\color{partcolor} Valoare} & {\sffamily\bfseries\color{partcolor} Efect} \\
\code{block}        & Element pe linie noua, ocupa toata latimea disponibila. Default pentru \code{<div>}, \code{<p>}, \code{<h1>...<h6>}, \code{<section>}. \\
\code{inline}       & Curge in textul din jur, ignora \code{width}/\code{height}. Default pentru \code{<span>}, \code{<a>}, \code{<strong>}, \code{<em>}. \\
\code{inline-block} & Curge in text ca inline, dar accepta \code{width}/\code{height} ca un block. Util pentru butoane in text. \\
\code{flex}         & Activeaza modelul Flexbox (uni-dimensional) pe \emph{copii}i elementului. Cel mai folosit pentru header-uri, navigatie, alinieri. \\
\code{grid}         & Activeaza modelul CSS Grid (bi-dimensional) pe copii. Pentru galerii, dashboard-uri, layout-uri complexe. \\
\code{none}         & Ascunde complet elementul (nu ocupa spatiu). Util pentru meniuri mobile inchise, modale ascunse. \\
\end{tabular}
\end{tcolorbox}

\begin{atentie}
\code{display: none} \textbf{ascunde complet} elementul (cititoarele de ecran il sar). Daca vrei doar sa-l faci invizibil dar accesibil, foloseste \code{visibility: hidden} sau \code{opacity: 0; pointer-events: none}.
\end{atentie}

% =========================================================
\section{Flexbox}

Flexbox (\emph{Flexible Box Layout}) este un model de layout uni-dimensional, util pentru distribuirea spatiului si alinierea continutului pe o axa (orizontala sau verticala).

\begin{lstlisting}[style=css, caption={Header cu logo si navigare aliniate}]
header {
  display: flex;                   /* activeaza flexbox */
  flex-direction: row;             /* axa principala = orizontala (default) */
  justify-content: space-between;  /* distribuire pe axa principala */
  align-items: center;             /* aliniere pe axa secundara */
  flex-wrap: wrap;                 /* trecere pe randul urmator daca nu incape */
  gap: var(--space-2);             /* spatiu intre copii */
}
\end{lstlisting}

\subsection{Cele 6 proprietati pe care le folosesti}

\begin{itemize}[itemsep=2pt]
  \item \code{display: flex} \textendash{} activeaza modelul (pe parinte).
  \item \code{flex-direction} \textendash{} axa principala: \code{row} (default, orizontal), \code{column} (vertical), \code{row-reverse}, \code{column-reverse}.
  \item \code{justify-content} \textendash{} distributie pe axa principala. Vezi tabelul de mai jos.
  \item \code{align-items} \textendash{} aliniere pe axa secundara. Vezi tabelul de mai jos.
  \item \code{gap} \textendash{} spatiu uniform intre copii (in unitati: \code{1rem}, \code{16px}). Inlocuieste \code{margin}.
  \item \code{flex-wrap} \textendash{} \code{nowrap} (default, totul pe o linie), \code{wrap} (trece pe linia urmatoare).
  \item \code{flex: 1} (pe copil) \textendash{} cre\textcommabelow{s}te pentru a umple spatiul disponibil.
\end{itemize}

\subsubsection{Valorile pentru \texttt{justify-content} (axa principala)}

\begin{tcolorbox}[
  enhanced, breakable, sharp corners, arc=0pt,
  colback=accentSoft, colframe=partcolor, boxrule=0pt, leftrule=3pt,
  left=12pt, right=12pt, top=8pt, bottom=8pt, fontupper=\small,
]
\renewcommand{\arraystretch}{1.35}
\begin{tabular}{@{}p{0.30\linewidth}@{\hspace{8pt}}p{0.66\linewidth}@{}}
\code{flex-start} (default) & Toate copii lipiti la inceputul axei. \\
\code{center}               & Grupate la mijlocul axei. \\
\code{flex-end}             & Lipiti la finalul axei. \\
\code{space-between}        & Primul/ultimul la margini, spa\textcommabelow{t}iu egal intre restul. \\
\code{space-around}         & Spa\textcommabelow{t}iu egal in jurul fiecaruia (margine = jumatate). \\
\code{space-evenly}         & Spa\textcommabelow{t}iu absolut egal peste tot (inclusiv la margini). \\
\end{tabular}
\end{tcolorbox}

\subsubsection{Valorile pentru \texttt{align-items} (axa secundara)}

\begin{tcolorbox}[
  enhanced, breakable, sharp corners, arc=0pt,
  colback=accentSoft, colframe=partcolor, boxrule=0pt, leftrule=3pt,
  left=12pt, right=12pt, top=8pt, bottom=8pt, fontupper=\small,
]
\renewcommand{\arraystretch}{1.35}
\begin{tabular}{@{}p{0.30\linewidth}@{\hspace{8pt}}p{0.66\linewidth}@{}}
\code{stretch} (default) & Copii se intind pe toata latimea axei secundare. \\
\code{center}            & Centrate pe axa secundara. \\
\code{flex-start}        & Lipiti la inceputul axei (sus pentru row, stanga pentru column). \\
\code{flex-end}          & Lipiti la final (jos / dreapta). \\
\code{baseline}          & Aliniati dupa linia de baza a textului (util pentru fonturi de marimi diferite). \\
\end{tabular}
\end{tcolorbox}

\begin{ponturi}
\textbf{Trick uzual:} \code{display: flex; justify-content: center; align-items: center;} centreaza orice element in containerul sau, atat orizontal cat si vertical, in 3 linii.
\end{ponturi}

% =========================================================
\section{CSS Grid (introducere scurta)}

Grid este sistemul de layout bidimensional al CSS-ului. Spre deosebire de flexbox (uni-dimensional), grid permite controlul simultan al randurilor si coloanelor.

\begin{lstlisting}[style=css, caption={Galerie cu coloane responsive automat}]
.gallery {
  display: grid;
  grid-template-columns: repeat(auto-fit, minmax(240px, 1fr));
  gap: var(--space-2);
}
\end{lstlisting}

Mecanism: \code{auto-fit} creeaza cate coloane incap; \code{minmax(240px, 1fr)} cere o latime minima de 240px si o expansiune egala (\code{1fr}) cand ramane spatiu. Rezultatul: galerie responsive fara media queries.

% =========================================================
\section{Mobile-first}

Strategia \emph{mobile-first} consta in scrierea regulilor pentru ecrane mici ca valori implicite, urmata de extinderea acestora pentru ecrane mai mari prin \emph{media queries}:

\begin{lstlisting}[style=css]
/* Stiluri implicite, pentru mobil */
.nav { display: none; }
.nav-toggle { display: flex; }

/* Adaugiri pentru ecrane >= 768px */
@media (min-width: 768px) {
  .nav { display: flex; }
  .nav-toggle { display: none; }
}
\end{lstlisting}

\begin{dano}[Asa nu \textendash{} desktop-first cu \texttt{max-width}]
\begin{lstlisting}[style=css, numbers=none]
.nav { display: flex; }
@media (max-width: 767px) {
  .nav { display: none; }
  .nav-toggle { display: flex; }
}
\end{lstlisting}
\textbf{Problema.} Mobilul descarca si proceseaza CSS conceput pentru desktop, apoi il anuleaza. Pierdere de performanta pe contextul cu cele mai limitate resurse.
\end{dano}

\begin{dado}[Asa da \textendash{} mobile-first cu \texttt{min-width}]
\begin{lstlisting}[style=css, numbers=none]
.nav { display: none; }
.nav-toggle { display: flex; }
@media (min-width: 768px) {
  .nav { display: flex; }
  .nav-toggle { display: none; }
}
\end{lstlisting}
\textbf{Beneficiu.} Mobilul (>60\% din trafic) primeste varianta optimizata implicit; desktopul adauga incremental.
\end{dado}

% =========================================================
\section{Suport pentru tema sistem (light / dark)}

Functia \code{light-dark()} si proprietatea \code{color-scheme} permit afisarea conditionata a culorilor in functie de preferinta sistemului de operare al utilizatorului:

\begin{lstlisting}[style=css]
html { color-scheme: light dark; }

:root {
  --bg:   light-dark(#FFFCF6, #0F172A);
  --text: light-dark(#1A1A1A, #F1F5F9);
}
\end{lstlisting}

Browserul selecteaza prima sau a doua valoare in functie de tema activa. Suportul nativ in browserele moderne este disponibil incepand cu 2024 (Chrome 123+, Firefox 120+, Safari 17.5+).

% =========================================================
\section{Pozi\textcommabelow{t}ionarea elementelor}

Proprietatea \code{position} controleaza modul in care un element se aseaza in fluxul paginii. Cinci valori, cu reper de cand le folose\textcommabelow{s}ti:

\begin{tcolorbox}[
  enhanced, breakable, sharp corners, arc=0pt,
  colback=accentSoft, colframe=partcolor, boxrule=0pt, leftrule=3pt,
  left=12pt, right=12pt, top=8pt, bottom=8pt, fontupper=\small,
]
\renewcommand{\arraystretch}{1.4}
\begin{tabular}{@{}p{0.20\linewidth}@{\hspace{6pt}}p{0.50\linewidth}@{\hspace{6pt}}p{0.24\linewidth}@{}}
{\sffamily\bfseries\color{partcolor} Valoare} & {\sffamily\bfseries\color{partcolor} Comportament} & {\sffamily\bfseries\color{partcolor} Cand foloseste} \\
\code{static} (default) & Respecta fluxul natural; \code{top}/\code{left} ignorate. & implicit, nu il setezi \\
\code{relative}         & Stay-in-place pana il mu\textcommabelow{t}i cu \code{top}/\code{left}; pastreaza spa\textcommabelow{t}iul. & ancora pentru \code{absolute} \\
\code{absolute}         & Scos din flux; pozi\textcommabelow{t}ionat fa\textcommabelow{t}a de primul parinte \code{relative}. & tooltip-uri, badge-uri \\
\code{fixed}            & Fixat pe viewport; ramane la scroll. & modal, buton ,,scroll-to-top'' \\
\code{sticky}           & \code{relative} pana ajunge la o margine, apoi \code{fixed}. & header sticky, sidebar \\
\end{tabular}
\end{tcolorbox}

\begin{atentie}
\code{absolute} are nevoie de un parinte \code{relative} (sau cade pana la viewport). Daca elementul ,,sare in lume'', verifica daca containerul are \code{position: relative}.
\end{atentie}

\begin{lstlisting}[style=css, caption={Header sticky modern}]
header {
  position: sticky;
  top: 0;
  background: var(--bg);
  z-index: 10;
}
\end{lstlisting}

% =========================================================
\section{Cheatsheet \textendash{} alte proprieta\textcommabelow{t}i cu enum-uri}

Cele patru cele mai uzuale dupa cele de mai sus. Le intalne\textcommabelow{s}ti zilnic.

\begin{tcolorbox}[
  enhanced, breakable, sharp corners, arc=0pt,
  colback=accentSoft, colframe=partcolor, boxrule=0pt, leftrule=3pt,
  left=12pt, right=12pt, top=8pt, bottom=8pt, fontupper=\small,
]
\renewcommand{\arraystretch}{1.35}
\begin{tabular}{@{}p{0.22\linewidth}@{\hspace{6pt}}p{0.34\linewidth}@{\hspace{6pt}}p{0.40\linewidth}@{}}
{\sffamily\bfseries\color{partcolor} Proprietate} & {\sffamily\bfseries\color{partcolor} Valori} & {\sffamily\bfseries\color{partcolor} Cand foloseste} \\
\code{box-sizing}     & \code{content-box} (default), \code{border-box} & \emph{Mereu border-box} \textendash{} aplicat global. \\
\code{overflow}       & \code{visible} (default), \code{hidden}, \code{scroll}, \code{auto} & \code{hidden} pentru a tunde, \code{auto} pentru scroll doar la nevoie. \\
\code{text-align}     & \code{left} (default), \code{center}, \code{right}, \code{justify} & Aliniere text in container. \\
\code{cursor}         & \code{auto}, \code{pointer}, \code{default}, \code{text}, \code{grab}, \code{not-allowed}, \code{wait} & \code{pointer} pe elemente clicabile non-link. \\
\code{font-weight}    & \code{normal} (400), \code{bold} (700), 100\textendash{}900 in pa\textcommabelow{s}i de 100 & \code{500} pentru \emph{medium}, \code{600} pentru \emph{semibold}. \\
\code{text-transform} & \code{none} (default), \code{uppercase}, \code{lowercase}, \code{capitalize} & In CSS, nu in HTML \textendash{} pastrezi semantica. \\
\code{visibility}     & \code{visible} (default), \code{hidden} & Ascunde dar pastreaza spa\textcommabelow{t}iul; \code{display: none} il scoate complet. \\
\code{pointer-events} & \code{auto} (default), \code{none} & \code{none} dezactiveaza click \textcommabelow{s}i hover (overlay-uri inactive). \\
\end{tabular}
\end{tcolorbox}

\begin{ponturi}
\textbf{Regula generala:} valori unitare ca \code{cursor: pointer}, \code{text-align: center}, \code{overflow: hidden} arata exact ca in tooltip-ul autocomplete al VS Code. \code{Ctrl+Space} dupa numele proprieta\textcommabelow{t}ii afi\textcommabelow{s}eaza lista completa.
\end{ponturi}

% =========================================================
\section{Anima\textcommabelow{t}ii CSS}

Animatiile imbunatatesc perceptia calitatii produsului \textcommabelow{s}i orienteaza aten\textcommabelow{t}ia utilizatorului. CSS suporta nativ doua mecanisme: \emph{transitions} pentru schimbari intre stari, si \emph{keyframes} pentru animatii arbitrare.

\subsection{Transitions \textendash{} schimbari intre stari}

\code{transition} interpoleaza fluent o proprietate cand valoarea sa se schimba (de exemplu, la \code{:hover}).

\begin{lstlisting}[style=css, caption={Buton cu hover smooth}]
.cta {
  background: var(--primary);
  color: white;
  padding: var(--space-1) var(--space-3);
  border-radius: var(--radius);
  transition: background 0.2s ease, transform 0.2s ease;
}

.cta:hover {
  background: var(--accent);
  transform: translateY(-2px);
}
\end{lstlisting}

Sintaxa: \code{transition: <proprietate> <durata> <timing-function> <delay>}.

\subsubsection{Timing function \textendash{} cum se simte miscarea}

Durata spune \emph{cat} dureaza animatia; \emph{timing function} spune \emph{cum} se distribuie miscarea in timp. Valorile predefinite sunt suficiente pentru 90\% din cazuri:

\begin{itemize}[itemsep=1pt]
  \item \code{linear} \textendash{} viteza constanta. Util pentru spinner-e si bare de progres.
  \item \code{ease} (default) \textendash{} accelereaza la inceput, decelereaza la final.
  \item \code{ease-in} \textendash{} porneste lent, se incheie rapid. Folosit pentru iesiri (\emph{exit}).
  \item \code{ease-out} \textendash{} porneste rapid, se incheie lent. Folosit pentru intrari (\emph{enter}); senzatie de aterizare lina.
  \item \code{ease-in-out} \textendash{} simetric. Folosit pentru tranzitii reversibile (toggle).
\end{itemize}

\begin{ponturi}[Regula UX rapida]
Pentru micro-interactiuni (hover, click): \code{ease-out 150ms\textendash{}250ms}. Sub 150ms se simte abrupt; peste 300ms se simte lent.
\end{ponturi}

\subsubsection{Cubic-bezier \textendash{} curba personalizata}

Cele 5 valori predefinite sunt scurtaturi pentru curbe Bezier de gradul 3. Cand vrei o curba personalizata (ex. \emph{bounce}, \emph{overshoot}), o defines\textcommabelow{t}i prin 4 numere care reprezinta cele doua puncte de control:

\begin{lstlisting}[style=css, caption={Easing personalizat cu overshoot subtil}]
.cta {
  transition: transform 0.3s cubic-bezier(0.34, 1.56, 0.64, 1);
}
.cta:hover { transform: scale(1.05); }
\end{lstlisting}

\begin{ponturi}
\textbf{Tool recomandat:} \url{https://cubic-bezier.com} \textendash{} editor vizual cu preview. Copiezi cele 4 valori si le lipesti in CSS.
\end{ponturi}

\subsection{Transform \textendash{} scalare, rota\textcommabelow{t}ie, transla\textcommabelow{t}ie}

\code{transform} aplica o transformare geometrica fara a afecta layout-ul (alte elemente nu se mut\u a). Operatiile uzuale:

\begin{itemize}
  \item \code{translate(x, y)} / \code{translateX} / \code{translateY} \textendash{} mutare pe axe (px, \%, rem). Mai performant decat \code{top}/\code{left}.
  \item \code{scale(n)} \textendash{} marire/micsorare proportionala. \code{scaleX}/\code{scaleY} pe o singura axa.
  \item \code{rotate(deg)} \textendash{} rotire in plan; pozitiv = sensul acelor de ceas.
  \item \code{skew(deg)} \textendash{} forfecare (rar folosit, util pentru efecte oblique).
  \item Combinare \emph{ordonata}: \code{transform: translateY(-2px) scale(1.05) rotate(2deg)} \textendash{} ordinea conteaza, se aplica de la dreapta la stanga.
\end{itemize}

\subsubsection{Transform-origin \textendash{} punctul de pivot}

Implicit, transformarile se aplica fata de centrul elementului. \code{transform-origin} schimba acest pivot:

\begin{lstlisting}[style=css, caption={Card care se ,,deschide'' din coltul stanga-sus}]
.card {
  transform-origin: top left;
  transition: transform 0.3s ease-out;
}
.card:hover { transform: scale(1.05); }
\end{lstlisting}

Valori uzuale: \code{center} (default), \code{top}, \code{bottom}, \code{left}, \code{right}, combinatii (\code{top left}, \code{50\% 0\%}), sau coordonate explicite in px/\%.

\begin{ponturi}[Performance]
Browserul accelereaza pe GPU doar \code{transform} si \code{opacity}. Anima\textcommabelow{t}ia altor proprieta\textcommabelow{t}i (ex. \code{width}, \code{margin}, \code{top}) declan\textcommabelow{s}eaza reflow \textcommabelow{s}i scade FPS. Regula: pentru mi\textcommabelow{s}care, folose\textcommabelow{s}te \code{transform: translate(...)}, nu \code{top}/\code{left}.
\end{ponturi}

\subsection{Keyframes \textendash{} animatii arbitrare}

Pentru anima\textcommabelow{t}ii care nu sunt declan\textcommabelow{s}ate de o stare (loading spinner, fade-in, parallax), \code{@keyframes} defineste o secven\textcommabelow{t}a de stari:

\begin{lstlisting}[style=css, caption={Spinner infinit}]
@keyframes rotate {
  from { transform: rotate(0deg); }
  to   { transform: rotate(360deg); }
}

.spinner {
  width: 32px;
  height: 32px;
  border: 4px solid #E5E7EB;
  border-top-color: var(--primary);
  border-radius: 50%;
  animation: rotate 1s linear infinite;
}
\end{lstlisting}

\begin{lstlisting}[style=css, caption={Fade-in pentru titlu hero}]
@keyframes fadeInUp {
  from { opacity: 0; transform: translateY(20px); }
  to   { opacity: 1; transform: translateY(0); }
}

.hero h1 {
  animation: fadeInUp 0.6s ease-out;
}
\end{lstlisting}

Sintaxa shorthand pentru animation: \code{animation: <nume> <durata> <timing> <delay> <iteratii> <direction> <fill-mode>}.

\subsubsection{Keyframes multi-step}

\code{@keyframes} accepta procente intermediare, nu doar \code{from} si \code{to}. Astfel po\textcommabelow{t}i descrie traiectorii complexe intr-o singura anima\textcommabelow{t}ie:

\begin{lstlisting}[style=css, caption={Buton care ,,pulseaz\u a'' atragator}]
@keyframes pulse {
  0%   { transform: scale(1);    box-shadow: 0 0 0 0 rgba(250, 165, 30, 0.7); }
  50%  { transform: scale(1.05); box-shadow: 0 0 0 12px rgba(250, 165, 30, 0); }
  100% { transform: scale(1);    box-shadow: 0 0 0 0 rgba(250, 165, 30, 0); }
}

.cta-featured { animation: pulse 2s ease-in-out infinite; }
\end{lstlisting}

\subsubsection{Animation-fill-mode \textendash{} ce ramane dupa anima\textcommabelow{t}ie}

Implicit, dupa ce animatia se termina, elementul revine la starea CSS originala. \code{animation-fill-mode} controleaza acest comportament:

\begin{itemize}[itemsep=1pt]
  \item \code{none} (default) \textendash{} elementul revine la stilul original.
  \item \code{forwards} \textendash{} pastreaza valorile din ultimul keyframe. \textbf{Cel mai folosit} pentru fade-in / slide-in (vrei sa ramana la \code{opacity: 1}).
  \item \code{backwards} \textendash{} aplica primul keyframe inainte de start (in timpul \code{animation-delay}).
  \item \code{both} \textendash{} combinatie \code{forwards + backwards}.
\end{itemize}

\begin{lstlisting}[style=css, caption={Fade-in care r\u amane vizibil}]
.hero h1 {
  opacity: 0;
  animation: fadeInUp 0.6s ease-out forwards;
}
\end{lstlisting}

\subsubsection{Animation-play-state \textendash{} pauz\u a \textcommabelow{s}i reluare}

\code{animation-play-state} permite punerea pe pauza a unei animatii \textendash{} util pentru carusele care se opresc la hover:

\begin{lstlisting}[style=css, numbers=none]
.carousel { animation: scroll 30s linear infinite; }
.carousel:hover { animation-play-state: paused; }
\end{lstlisting}

\subsection{Pattern complet: card cu hover-lift}

Combinatia transition + transform + box-shadow produce efectul de ,,ridicare'' pe care il vezi pe orice landing page modern:

\begin{lstlisting}[style=css, caption={Hover-lift, pattern de baz\u a pentru carduri}]
.card {
  background: var(--bg);
  padding: var(--space-3);
  border-radius: var(--radius);
  box-shadow: 0 1px 3px rgba(0, 0, 0, 0.08);
  transition: transform 0.25s ease-out, box-shadow 0.25s ease-out;
  cursor: pointer;
}

.card:hover {
  transform: translateY(-4px);
  box-shadow: 0 12px 24px rgba(0, 0, 0, 0.12);
}
\end{lstlisting}

Detalii care fac diferenta:
\begin{itemize}[itemsep=1pt]
  \item \code{ease-out} pe enter \textendash{} senzatie de aterizare lina.
  \item Umbra cre\textcommabelow{s}te si se ,,rasfira'' (offset Y mai mare, blur mai mare) \textendash{} simuleaza distan\textcommabelow{t}a fa\textcommabelow{t}a de fundal.
  \item Tranzi\textcommabelow{t}ie pe ambele proprieta\textcommabelow{t}i simultan, nu doar \code{transform}.
  \item Cursor \code{pointer} \textendash{} afi\textcommabelow{s}eaza inten\textcommabelow{t}ia ca elementul e clicabil.
\end{itemize}

\subsection{Accesibilitate: \texttt{prefers-reduced-motion}}

Unii utilizatori (afec\textcommabelow{t}iuni vestibulare, dependente de aten\textcommabelow{t}ie) au setat in OS preferin\textcommabelow{t}a pentru anima\textcommabelow{t}ii reduse. CSS detecteaza aceasta cerere:

\begin{lstlisting}[style=css]
@media (prefers-reduced-motion: reduce) {
  *, *::before, *::after {
    animation-duration: 0.01ms !important;
    transition-duration: 0.01ms !important;
  }
}
\end{lstlisting}

Conformitatea este o cerin\textcommabelow{t}a WCAG 2.1 (criteriul 2.3.3) \textcommabelow{s}i o practica recomandata pentru orice site cu animatii.

% =========================================================
\begin{bonusbox}
\textbf{Acest box e optional.} Daca esti la inceput, treci direct la \emph{Verificare cuno\textcommabelow{s}tin\textcommabelow{t}e}. Sec\textcommabelow{t}iunile de mai jos sunt pentru cand vrei sa ridici stacheta.

\medskip

\textbf{\faCube\ Transform\u ari 3D \textendash{} flip-card}

CSS suporta nativ transformari in spatiu, nu doar in plan. Trei piese:

\begin{itemize}[itemsep=1pt]
  \item \code{perspective} (pe parinte) \textendash{} ,,distan\textcommabelow{t}a'' fata de plan. Valori uzuale 600\textendash{}1500px.
  \item \code{transform-style: preserve-3d} (pe parinte) \textendash{} permite copiilor sa existe in spatiu 3D.
  \item \code{backface-visibility: hidden} (pe fete) \textendash{} ascunde fa\textcommabelow{t}a din spate.
\end{itemize}

\begin{lstlisting}[style=css, caption={Card care se intoarce la hover (flip-card)}]
.flip-card { perspective: 1000px; }
.flip-card-inner {
  position: relative;
  width: 100%; height: 100%;
  transition: transform 0.6s;
  transform-style: preserve-3d;
}
.flip-card:hover .flip-card-inner { transform: rotateY(180deg); }
.flip-card-front, .flip-card-back {
  position: absolute; inset: 0;
  backface-visibility: hidden;
}
.flip-card-back { transform: rotateY(180deg); }
\end{lstlisting}

\medskip

\textbf{\faScroll\ Anima\textcommabelow{t}ii legate de scroll \textendash{} \texttt{animation-timeline}}

CSS modern (Chrome 115+, 2023) permite legarea unei animatii de pozitia scroll-ului, fara JavaScript:

\begin{lstlisting}[style=css, caption={Bara de progres care urmare\textcommabelow{s}te scroll-ul}]
@keyframes growProgress {
  from { transform: scaleX(0); }
  to   { transform: scaleX(1); }
}

.progress-bar {
  position: fixed; top: 0; left: 0;
  height: 4px; width: 100%;
  background: var(--accent);
  transform-origin: left;
  animation: growProgress linear;
  animation-timeline: scroll(root);  /* legat de scroll-ul paginii */
}
\end{lstlisting}

\code{view-timeline} este varianta complementar\u a: anima\textcommabelow{t}ia avanseaza pe masura ce un element intra in viewport. Util pentru fade-in la scroll, dar fara cost JavaScript.

\begin{nota}
Suportul pentru \code{animation-timeline} este partial in 2026: Chrome/Edge da, Firefox in spatele unui flag, Safari in pregatire. Foloseste cu fallback (animatia pur si simplu nu ruleaza in browserele vechi \textendash{} progressive enhancement OK).
\end{nota}

\medskip

\textbf{\faMagic\ \texttt{filter} \textcommabelow{s}i \texttt{backdrop-filter} \textendash{} efecte de imagine}

\code{filter} aplica un efect vizual pe element. \code{backdrop-filter} aplica acela\textcommabelow{s}i efect pe ce este \emph{in spate}: a\textcommabelow{s}a se face glassmorphism-ul (efect ,,sticla mata'').

\begin{lstlisting}[style=css, caption={Header transparent cu blur pe continutul dedesubt}]
header {
  background: rgba(255, 252, 246, 0.7);
  backdrop-filter: blur(12px) saturate(180%);
  -webkit-backdrop-filter: blur(12px) saturate(180%);  /* Safari */
}
\end{lstlisting}

Functii utile: \code{blur(8px)}, \code{brightness(1.2)}, \code{contrast(1.1)}, \code{grayscale(80\%)}, \code{hue-rotate(15deg)}, \code{drop-shadow(0 4px 8px rgba(0,0,0,0.2))}. Pot fi combinate prin spatiere.

\medskip

\textbf{\faCut\ \texttt{clip-path} \textendash{} forme custom \textcommabelow{s}i reveal effects}

\code{clip-path} ,,decupeaza'' elementul dupa o forma. Util pentru hero-uri oblice, butoane in forma de chevron, anima\textcommabelow{t}ii de tip ,,revelare'':

\begin{lstlisting}[style=css, caption={Imagine cu margine diagonala}]
.hero-image {
  clip-path: polygon(0 0, 100% 0, 100% 85%, 0 100%);
}
\end{lstlisting}

\begin{lstlisting}[style=css, caption={Reveal animat: pagina se ,,deschide'' la load}]
@keyframes reveal {
  from { clip-path: inset(0 100% 0 0); }
  to   { clip-path: inset(0 0 0 0); }
}
.hero h1 { animation: reveal 0.8s ease-out forwards; }
\end{lstlisting}

\begin{ponturi}
\textbf{Tool vizual:} \url{https://bennettfeely.com/clippy/} \textendash{} editor SVG-like care genereaza valoarea \code{clip-path} pentru orice forma desenata.
\end{ponturi}

\medskip

\textbf{\faMobile*\ Container queries \textendash{} ,,media queries'' pentru componente}

Problema cu \code{@media}: depinde de viewport, nu de containerul real al componentei. Un card de \code{300px} se comporta diferit intr-o coloana ingusta vs pe ecran lat, dar \code{@media} nu vede asta.

\code{@container} (CSS 2023, suport stabil) rezolva exact asta:

\begin{lstlisting}[style=css, caption={Card responsive in functie de container, nu de viewport}]
.card-wrapper { container-type: inline-size; }

.card { display: block; }

@container (min-width: 400px) {
  .card { display: flex; gap: var(--space-2); }
}
\end{lstlisting}

Acela\textcommabelow{s}i card devine ,,horizontal'' cand are spatiu, ,,vertical'' cand nu \textendash{} fara sa-i pese de marimea ecranului. Standardul modern pentru design system-uri.

\medskip

\textbf{\faRoute\ Unde mergi de aici}

\begin{itemize}[itemsep=1pt]
  \item \emph{Animatable properties:} \url{https://developer.mozilla.org/en-US/docs/Web/CSS/CSS_animated_properties}
  \item \emph{Scroll-driven animations:} \url{https://developer.chrome.com/docs/css-ui/scroll-driven-animations}
  \item \emph{Container queries:} \url{https://web.dev/learn/design/container-queries}
\end{itemize}
\end{bonusbox}

% =========================================================
\section*{Verificare cuno\textcommabelow{s}tin\textcommabelow{t}e}

\begin{selfcheck}
\begin{enumerate}[itemsep=2pt]
  \item Care este specificitatea pentru selectorul \code{.btn.primary}? Comparativ, ce specificitate are \code{button\#submit}?
  \item De ce strategia mobile-first este preferabila pentru contextul actual al traficului web?
  \item \emph{Spot the issue:} \code{.btn \{ background: red !important; \}} repetat in trei locuri.
  \item Care este diferen\textcommabelow{t}a func\textcommabelow{t}ionala intre \code{rem} si \code{px} pentru \code{font-size}?
  \item Ce face \code{box-sizing: border-box} fa\textcommabelow{t}a de comportamentul implicit?
  \item Pentru un hover smooth pe un buton, ce duration si timing-function alegi? Argumenteaza.
  \item De ce \code{transform: translateY(-2px)} pe hover este preferabil fa\textcommabelow{t}a de \code{top: -2px}?
  \item Cand folose\textcommabelow{s}ti \code{animation-fill-mode: forwards}? Da un exemplu concret.
\end{enumerate}
\end{selfcheck}

% =========================================================
\section*{Recapitulare}

\begin{recap}[Faza 2 -- CSS]
\begin{itemize}[itemsep=1pt]
  \item Stiluri externe in \code{<link rel="stylesheet">}; nu inline.
  \item Specificitate: \code{(id, clasa, tag)}; \code{!important} este \emph{antipattern}.
  \item \code{box-sizing: border-box} aplicat global pentru calcul predictibil.
  \item Unitati: \code{rem} pentru fonturi/spatiere; \code{px} pentru border; \code{\%}/\code{vw}/\code{vh} pentru layout responsiv.
  \item Variabilele CSS in \code{:root}; redefinire locala pentru teme contextuale.
  \item Flexbox: 6 proprietati cheie (\code{display}, \code{flex-direction}, \code{justify-content}, \code{align-items}, \code{gap}, \code{flex: 1}).
  \item Mobile-first cu \code{@media (min-width: ...)}, nu desktop-first cu \code{max-width}.
  \item Pozitionare: \code{static} (default), \code{relative}, \code{absolute}, \code{fixed}, \code{sticky}.
  \item Anima\textcommabelow{t}ii: \code{transition} pentru schimbari de stare (hover/focus), \code{@keyframes} pentru animatii arbitrare.
  \item Timing function: \code{ease-out} pentru intrari, \code{ease-in} pentru iesiri, \code{cubic-bezier()} pentru curbe custom.
  \item Performance: anima doar \code{transform} \textcommabelow{s}i \code{opacity} (accelerate GPU). Evita \code{top}/\code{left}/\code{width} \textendash{} declan\textcommabelow{s}eaza reflow.
  \item Pattern hover-lift: \code{transform: translateY(-4px) + box-shadow} cu \code{transition} pe ambele.
  \item \code{animation-fill-mode: forwards} pastreaza ultimul keyframe (esential pentru fade-in care ramane vizibil).
  \item Accesibilitate: respecta \code{prefers-reduced-motion}.
  \item \faRocket\ Bonus: 3D transforms, scroll-driven animations, \code{filter}, \code{clip-path}, container queries.
\end{itemize}
\end{recap}

\partdivider

% =========================================================
\begin{joc}
\textbf{Platforma:} Flexbox Froggy (\url{https://flexboxfroggy.com}).

\medskip

\textbf{Activitatea.} Aplicatia este un set de 24 de provocari care necesita pozitionarea unor obiecte pe o grila prin folosirea proprietatilor \code{justify-content}, \code{align-items}, \code{flex-direction}, \code{order}, \code{flex-wrap}.

\medskip

\textbf{Obiectiv:} primele 6 niveluri (acopera \code{justify-content} si \code{align-items}, cele mai folosite proprietati flexbox in practica).

\medskip

\textbf{Durata estimata:} 5 minute.
\end{joc}

% =========================================================
\section*{\faGraduationCap\ \ Resurse pentru aprofundare}

\begin{itemize}[itemsep=2pt]
  \item Mozilla Developer Network. \emph{CSS reference}. \url{https://developer.mozilla.org/en-US/docs/Web/CSS}
  \item Google. \emph{Learn CSS}. \url{https://web.dev/learn/css}
  \item Coyier C. \emph{A Complete Guide to Flexbox}. CSS-Tricks. \\\url{https://css-tricks.com/snippets/css/a-guide-to-flexbox/}
  \item Coyier C. \emph{A Complete Guide to Grid}. CSS-Tricks. \\\url{https://css-tricks.com/snippets/css/complete-guide-grid/}
  \item The Odin Project. \emph{CSS Foundations}. \\\url{https://www.theodinproject.com/paths/foundations/courses/foundations\#css-foundations}
  \item Aplicatii similare: \url{https://flukeout.github.io} (selectori), \url{https://cssgridgarden.com} (grid), \url{https://cssbattle.dev} (speedrun).
\end{itemize}

\bigskip

\bridge{Capitolul 3 (JavaScript) adauga interactivitate. Pagina stilizata in Faza 2 va raspunde la click-uri, va schimba dark mode, va procesa formulare \textendash{} fara reincarcare.}
